\documentclass[10pt,a4paper]{amsart}
\usepackage[margin=19mm]{geometry}
\usepackage{amsmath,amssymb,mathtools}
\usepackage[scr=rsfs,cal=euler]{mathalfa}
\usepackage{xfrac}
\usepackage[unicode,hidelinks,bookmarks=false]{hyperref}
\newcommand{\ie}{i.e.\ }
\newcommand{\cf}{cf.\ }
\newcommand{\resp}{respectively\ }
\newcommand{\Subsection}{\subsection}
\providecommand{\nobreakdash}{\nobreak\hbox{-}}
\setlength{\emergencystretch}{2em}
\multlinegap0pt
\usepackage{bm}
\usepackage{bbm}
\usepackage{dsfont}
\usepackage{xspace}
\usepackage{cancel}
\usepackage{mathtools}
\usepackage{amsthm}
\makeatletter
\@ifundefined{theorem}{\newtheorem{theorem}{Theorem}}{}
\@ifundefined{claim}{\newtheorem{claim}[theorem]{Claim}}{}
\@ifundefined{lemma}{\newtheorem{lemma}[theorem]{Lemma}}{}
\makeatother
\title[Finding the diameter of a tree with distance queries]{Finding the diameter of a tree with distance queries}
\author{D\'aniel Gerbner, Andr\'as Imolay, Kartal Nagy, Bal\'azs Patk\'os, Krist\'of Z\'olomy}
\date{Source: arXiv:2509.23326v1 (CC BY 4.0)}
\begin{document}
\maketitle

\noindent\emph{Source and licence.} Finding the diameter of a tree with distance queries. Version arXiv:2509.23326v1 (CC BY 4.0), distributed under the Creative Commons Attribution 4.0 International licence, as recorded in the arXiv licence metadata for this exact version and shown on the abstract page https://arxiv.org/abs/2509.23326v1. Full rights notice: licenses/arxiv-2509.23326-CC-BY-4.0.txt.\par\smallskip

\noindent\textit{Editorial scope.} Counted unit: the Lemma whose source label is \texttt{lem:isom} in the article (source0.tex lines 650--652, source label \texttt{lem:isom}), delivered together with the article's own complete original proof of it (source0.tex lines 654--699, one \texttt{proof} environment of 46 lines). No mathematics is written, repaired or re-derived: statement and proof are the article's own text line for line. Required context retained uncounted: lemnew2 (lines 622--624). Every definition, display and label of those lines is retained unchanged. Omitted from this package and not redistributed: 1--621, 625--649, 653--653, 700--807. Nothing inside a retained excerpt is omitted or altered: every retained body is delivered exactly as the article's lines are written, so no omission receipt of any kind accompanies this package. Citations: the retained excerpts carry no citation command at all, so no bibliography entry of the article is transcribed and no reference is redistributed. Reference adaptation: none was needed and none was made; every reference inside the delivered unit resolves inside it, so no unresolved reference or citation is produced. Printed slips kept literal: no printed word, symbol, hypothesis or number of the counted statement or of its proof was altered. Layout and class: the article's own class is \texttt{amsart} with packages including amssymb, amsmath, amsthm, graphicx, float, hyperref, todonotes, soul; the wrapper is the lane's pinned \texttt{amsart} editorial wrapper at 10pt on A4, which restates the theorem-like environments the retained text needs and adds the article's own macro definitions that the retained text needs (none), with extra wrapper packages (bm, bbm, dsfont, xspace, cancel, mathtools). The delivered numbering is the unit's own. Assets: no retained span contains a figure environment, a graphics-inclusion command, a PDF-page inclusion, a table or a TikZ picture; no image file is copied into this package, the package declares no asset dependency and no image file is redistributed with it. Licence: version arXiv:2509.23326v1 of this article is distributed under the Creative Commons Attribution 4.0 International licence, as recorded in the arXiv licence metadata for this exact version and shown on the abstract page https://arxiv.org/abs/2509.23326v1. A full rights notice is delivered at licenses/arxiv-2509.23326-CC-BY-4.0.txt. Characters: every retained excerpt is pure ASCII, so the delivered engine, XeLaTeX with its default Latin Modern text font, reports no missing character.
\begin{lemma}\label{lemnew2}
    For every pair of consistent trees $T_0, T_1$, there is a vertex $z$ such that the distance of $z$ and $v$ is the same in $T_0$ and $T_1$ for every vertex $v$.
\end{lemma}

\begin{lemma}\label{lem:isom}
    For every pair of consistent trees $T_0, T_1$, there is a graph isomorphism $\varphi \colon V(T_0) \to V(T_1)$. Furthermore, if vertices $x$ and $y$ have distance at least 5 in $T_0$, then they have the same distance in $T_1$.
\end{lemma}

\begin{proof}
Let $u$ be a vertex such that the distance of $u$ and any other vertex $v$ is the same in $T_0$ and $T_1$. Recall $u$ exists by Lemma~\ref{lemnew2}. Let $U_\ell$ denote the set of vertices at distance $\ell$ from $u$ in $T_0$ and $T_1$. We consider $u$ to be the root of trees $T_0$ and $T_1$. In particular, we apply the usual terminology. For a tree $T \in \{T_0, T_1\}$ and a vertex $v\in U_\ell$, the \textit{parent} of $v$ is its neighbor in $U_{\ell-1}$ and its \textit{children} are its neighbors in $U_{\ell+1}$. The \textit{ancestors} of $v$ are the parent of $v$ and parents of its ancestors, the \textit{descendants} of $v$ are the children of $v$ and children of its descendants. The following trivial claim will be important later.

\begin{claim}\label{cla1}
  \textbf{(i)}  If $v\in U_\ell$ and $v'\in U_j$, then $v'$ is a descendant of $v$ in $T \in \{T_0, T_1\}$ if and only if their distance is $j-\ell$ in $T$.

  \textbf{(ii)} If $v\in U_\ell$ and $v'\in U_j$ is a descendant of $v$ in $T \in \{T_0, T_1\}$, then a vertex $v''\in U_i$ is a descendant of $v$ in $T$ if and only if the distance of $v'$ and $v''$ is at most $j+i-2\ell$ in $T$. 
\end{claim}

First, we prove the following.

\begin{claim}\label{cla2}
    If a vertex $v$ has at least three descendants in $T_0$, then the same vertices are its descendants in $T_1$. Similarly, if a vertex $v$ has at least three descendants in $T_1$, then the same vertices are its descendants in $T_0$.
\end{claim}

\begin{proof}[Proof of Claim~\ref{cla2}] 
At least one of the descendants $v'$ of $T_0$ was queried with $v$, thus we know that $v'$ is a descendant of $v$ in $T_1$ as well by Claim~\ref{cla1} (i).

If there is a vertex $v''$ that is a descendant of $v$ in $T_0$ and $v''$ is not a descendant of $v$ in $T_1$, then $v''$ has an ancestor $w\in U_\ell$ in $T_1$ with $w\neq v$. Then $v''$ was not queried with $v$ and $w$, thus $v''$ was queried with $v'$.
Therefore, we get that $v''$ is a descendant of $v$ in $T_1$ as well by Claim~\ref{cla1} (ii), a contradiction.
\end{proof}


We are ready to give an isomorphism $\varphi$ between $T_0$ and $T_1$. First, we define a subtree of $T_0$, where $\varphi$ will be the identity. We call the vertices of this subtree \emph{nice}. We define the nice vertices recursively with respect to the distance from $u$. Let $u$ be nice. Assume that we determined the nice vertices up to distance $\ell-1$ from $u$. We say that $v \in U_\ell$ is nice if the parent of $v$ in $T_0$ (and hence in $T_1$) is nice and the children of $v$ in $T_0$ and $T_1$ are the same. Clearly, the subgraph spanned by the nice vertices is a subtree of $T_0$ and $T_1$, and the identity map between them is an isomorphism of these subtrees. 
The goal is to extend this isomorphism to $V(T_0)$. 

Let $u'\in U_{\ell-1}$ be a nice vertex. It follows that the set of descendants of $u'$ in $T_0$ and $T_1$ are the same. Indeed, let $v$ be one of the closest descendants to $u'$ in $T_0$ that is not a descendant of $u'$ in $T_1$. As $u'$ is nice, $v$ cannot be a child of $u'$. Also, $v$ has distinct parents in $T_0$ and $T_1$, thus was not queried with them. But then $v$ was queried with $u'$, thus, by Claim \ref{cla1} (i), it is a descendant of $u'$ in both $T_0$ and $T_1$, a contradiction. 

Let $U_{\ell}'$ denote the set of not nice vertices in $U_\ell$ that are children of $u'$, i.e., whose descendants are not the same in $T_0$ and $T_1$. Now we will define $\varphi$ on $U'_\ell$ and the descendants of vertices in $U'_\ell$. By the earlier observations, these are the same set of vertices in $T_0$ and $T_1$. By Claim \ref{cla2}, the vertices in $U_{\ell}'$ have at most two descendants in both $T_0$ and $T_1$.
Consider first $u_0\in U_\ell'$ with two children $v,v'$ in $T_0$ such that the parent of $v$ is another vertex $u_1$ in $T_1$. Then $v$ was not queried with $u_0$, nor with $u_1$, thus we know $q(v,v')=2$, hence we know that $v'$ is a child of $u_1$ in $T_1$. By Claim \ref{cla2}, $u_0$ has no other descendants in $T_0$ and $u_1$ has no other descendants in $T_1$. Define $\varphi(v)=v$, $\varphi(v')=v'$ and $\varphi(u_0)=u_1$.
Conversely, if a vertex $u_1 \in U_{\ell}'$ has two children $v$ and $v'$ in $T_1$, then there is a vertex $u_0 \in U_{\ell}'$ with children $v, v'$ in $T_0$, hence we already defined the preimage of $u_1, v, v'$ with respect to $\varphi$ in this step.

Let $U_\ell''$ denote the set of vertices in $U_\ell'$ that do not have two children in $T_0$. As they have at most two descendants, we can partition $U_\ell''$ into $3$ sets $U_{\ell 0}''$, $U_{\ell 1}''$, $U_{\ell 2}''$ with $U_{\ell 0}''$ consisting of vertices with no children, $U_{\ell 1}''$ being the set of vertices that have one child with no children, and $U_{\ell 2}''$ having vertices with one child who has one child with no children $T_0$. Similarly, let $V_\ell''$ denote the set of vertices in $U_\ell'$ that do not have two children in $T_1$. We can partition $V_\ell''$ into $3$ sets $V_{\ell 0}''$, $V_{\ell 1}''$, $V_{\ell 2}''$ with $V_{\ell 0}''$ consisting of vertices with no children, $V_{\ell 1}''$ being the set of vertices that have one child with no children, and $V_{\ell 2}''$ having vertices with one child who has one child with no children in $T_1$. 
We have $|U_{\ell i}''|=|V_{\ell i}''|$ for $i=0,1,2$, as the number of descendants of $u'$ in $U_{\ell+1}$ and $U_{\ell+2}$ is the same in $T_0$ and $T_1$.
Define $\varphi$ on $U_\ell''$ by taking arbitrary bijections between vertices in $U_{\ell i}''$ and $V_{\ell i}''$ for $i=1,2,3$. Finally, if a vertex $x$ is the child of some vertex $z \in U_\ell''$, then let $\varphi(x)$ be the child of $\varphi(z)$, and similarly, if $x$ is the child of the child of some vertex $z \in U_\ell''$, then let $\varphi(x)$ be the child of the child of $\varphi(z)$.

Doing this for all nice vertices $u'$ defines the mapping $\varphi$ on every vertex of $V(T_0)$, and $\varphi$ is clearly an isomorphism by definition.

From the above definitions and observations, it is not hard to prove the second part of the lemma. Consider vertices $x$ and $y$ with distance at least 5 in $T_0$. Clearly, if $x$ or $y$ is nice, then their distance is the same in $T_0$ and $T_1$. Let $x'$ and $y'$ be the closest nice ancestors of $x$ and $y$, respectively. By the nice property, these are the same vertices in $T_0$ and $T_1$. It is also clear that in the case of $x' \neq y'$, the distance of $x$ and $y$ is the same in $T_0$ and $T_1$. Hence, assume that $x'=y'$ and let $x' \in U_{\ell-1}$. Let $x \in U_{\ell+a}$ and $y \in U_{\ell+b}$. As by Claim \ref{cla2}, children of $x'$ that are ancestors of either $x$ or $y$ have at most two descendants, we know that $a,b \leq 2$. 

Furthermore, as the distance of $x,y$  is at least 5 in $T_0$, $a+b \geq 3$ and $x$ and $y$ has $x'$ as the closest common ancestor in $T_0$. By Claim \ref{cla2}, there cannot be a common ancestor $u \in U_{l} \cup U_{l+1} \cup U_{l+2}$ with at least 3 descendants in $T_1$.

Assume first that $a=b=2$. Then $d(x,y)$ in $T_1$ could be either $2,4$, or $6$. If either $d(x,y)=2$ or $d(x,y)=4$, the common grandparent $u$ of $x$ and $y$ in $T_1$ has at least $3$ descendants, a contradiction by the previous observation.

Now, assume by symmetry that $a=2, b=1$. We need to show that $d(x,y)$ in $T_1$ cannot be $1$ or $3$. If $d(x,y)=3$, then their closest common ancestor in $T_1$ has at least $3$ descendants, which is again a contradiction. In the remaining case, $y$ is the parent of $x$ in $T_1$. Let $w_1$ be the parent of $x$, $w_2$ be the grandparent of $x$, and $w_3$ be the parent of $y$ in $T_0$. As we have two vertices, $y$ and $w_1$, whose distance from $x$ differs in $T_0$ and $T_1$, we must have that $w_2$ is the grandparent of $x$ in $T_1$ as well, hence it is the parent of $y$. Then the distances of $y$ from $x$, $w_2$, and $w_3$ are all different in $T_0$ and $T_1$, a contradiction.
\end{proof}

\end{document}
